% ===== PRÉAMBULE CANONIQUE — à coller VERBATIM en tête de CHAQUE .tex =====
\documentclass[11pt,a4paper]{article}
\usepackage[utf8]{inputenc}\usepackage[T1]{fontenc}\usepackage{lmodern}
\usepackage[french]{babel}
\usepackage{amsmath,amssymb,mathtools}\usepackage{icomma}
\usepackage{siunitx}\sisetup{locale=FR}  % \qty{}{}, \si{}, \ang{}
\usepackage{xcolor}\usepackage{colortbl}
\usepackage{tikz}\usepackage{pgfplots}\pgfplotsset{compat=1.18}
\usepackage[siunitx,europeanresistors]{circuitikz} % circuits (élec.)
\usepackage[most]{tcolorbox}\usepackage{enumitem}\usepackage{array,booktabs}
\usepackage[a4paper,margin=2.2cm]{geometry}
\usetikzlibrary{arrows.meta,calc,angles,quotes,positioning,patterns,%
  backgrounds,shapes.geometric,decorations.pathmorphing,decorations.markings,intersections}
\definecolor{primary}{RGB}{222,49,99}\definecolor{primarydark}{RGB}{150,22,55}
\definecolor{defcol}{RGB}{0,114,178}\definecolor{methcol}{RGB}{52,92,148}
\definecolor{excol}{RGB}{0,135,90}\definecolor{attncol}{RGB}{200,40,55}
\tcbset{boxbase/.style={enhanced,breakable,boxrule=0.9pt,arc=2.5pt,
  left=3mm,right=3mm,top=2.3mm,bottom=2.3mm,fonttitle=\bfseries,coltitle=white}}
% --- boîtes TUTORIEL (noms EXACTS, ne pas renommer) ---
\newtcolorbox{cle}[1][]{boxbase,colback=defcol!6,colframe=defcol,
  title={\textsc{À retenir}\ifx\relax#1\relax\relax\else\textmd{ : \itshape #1}\fi}}
\newtcolorbox{methode}[1][]{boxbase,colback=methcol!7,colframe=methcol,sharp corners,
  title={\textsc{Méthode}\ifx\relax#1\relax\relax\else\textmd{ --- \itshape #1}\fi}}
\newtcolorbox{exemplebox}[1][]{boxbase,colback=excol!5,colframe=excol,
  title={\textsc{Exemple}\ifx\relax#1\relax\relax\else\textmd{ : \itshape #1}\fi}}
\newtcolorbox{piege}[1][]{boxbase,colback=attncol!6,colframe=attncol,
  title={\textsc{Piège}\ifx\relax#1\relax\relax\else\textmd{ : \itshape #1}\fi}}
\newtcolorbox{remarque}[1][]{boxbase,colback=black!4,colframe=black!45,
  title={\textsc{Remarque}\ifx\relax#1\relax\relax\else\textmd{ : \itshape #1}\fi}}
% --- boîtes EXERCICE / CORRIGÉ (noms EXACTS, auto-comptées) ---
\newtcolorbox[auto counter]{exo}[1][]{boxbase,colback=primary!4,colframe=primary,
  title={\textsc{Exercice}~\thetcbcounter\ifx\relax#1\relax\relax\else\textmd{ --- \itshape #1}\fi}}
\newtcolorbox[auto counter]{corrige}[1][]{boxbase,colback=excol!4,colframe=excol,
  title={\textsc{Corrigé}~\thetcbcounter\ifx\relax#1\relax\relax\else\textmd{ --- \itshape #1}\fi}}
\newcommand{\vv}[1]{\vec{#1}}\newcommand{\ds}{\displaystyle}
\newcommand{\R}{\mathbb{R}}\newcommand{\N}{\mathbb{N}}\newcommand{\Z}{\mathbb{Z}}
% =========================================================================
\begin{document}
\sisetup{per-mode=symbol}
% ================= DIFFICILE =================

\begin{exo}[le satellite géostationnaire --- vrai ou faux]
Un satellite est géostationnaire (période égale à celle de rotation de la Terre,
orbite équatoriale, même sens). Indique si chaque affirmation est \textbf{vraie} ou
\textbf{fausse}.
\begin{enumerate}[label=\alph*)]
  \item Un observateur au sol le voit \emph{immobile}.
  \item Il a la même vitesse \emph{angulaire} $\omega$ qu'une maison située sur l'équateur.
  \item Il a la même vitesse \emph{linéaire} $v$ qu'une maison située sur l'équateur.
  \item Il existe plusieurs altitudes géostationnaires possibles.
\end{enumerate}
\end{exo}

\begin{exo}[classer trois satellites : $\omega$ et $v$]
Trois satellites sont aux altitudes $h_1>h_2>h_3$ (rayons $r_1>r_2>r_3$).
\begin{enumerate}[label=\alph*)]
  \item Sachant que $\omega=\dfrac{v}{r}=\sqrt{\dfrac{\mathcal{G}M_T}{r^3}}$, classe les
        vitesses \emph{angulaires} $\omega_1,\omega_2,\omega_3$.
  \item Avec $v=\sqrt{\dfrac{\mathcal{G}M_T}{r}}$, classe les vitesses \emph{linéaires}
        $v_1,v_2,v_3$.
\end{enumerate}
\end{exo}

\begin{exo}[à quelle altitude est-il géostationnaire ?]
On cherche l'altitude géostationnaire à partir de la période $T=\qty{86164}{\second}$
(un jour). On donne $\mathcal{G}M_T=\qty{4.0e14}{}$ (SI), $4\pi^2\approx 39{,}5$,
$R_T=\qty{6400}{\kilo\metre}$.
\begin{enumerate}[label=\alph*)]
  \item À partir de $T^2=\dfrac{4\pi^2 r^3}{\mathcal{G}M_T}$, isole $r$ et calcule-le
        (on utilisera $(\num{86164})^2\approx\num{7.42e9}$).
  \item Déduis-en l'altitude $h=r-R_T$.
\end{enumerate}
\end{exo}

\begin{exo}[même $\omega$, mais quelle vitesse ?]
Un satellite géostationnaire est à $r=\qty{4.22e7}{\metre}$ et une maison sur
l'équateur à $R_T=\qty{6.4e6}{\metre}$. Tous deux tournent à la même vitesse
angulaire $\omega=\qty{7.27e-5}{\radian\per\second}$.
\begin{enumerate}[label=\alph*)]
  \item Calcule la vitesse linéaire du satellite, $v_{\text{sat}}=\omega\,r$.
  \item Calcule celle de la maison, $v_{\text{maison}}=\omega\,R_T$.
  \item Combien de fois le satellite va-t-il plus vite que la maison ? Conclus sur le
        piège « même $\omega$ $\ne$ même $v$ ».
\end{enumerate}
\end{exo}

\end{document}
